%! Solving Polynomial Equations by Factoring — Unit 1, Lesson 1.6 — Exit Ticket
\documentclass[10pt]{article}
\usepackage{atda-article}
\usepackage{atda-boxes}
\newcommand{\TallMath}[1]{$\displaystyle #1\rule[-1.4em]{0pt}{3.2em}$}
\setlength{\workrowsep}{10pt}

\begin{document}

\pageheader{Unit 1, Lesson 1.6}{Exit Ticket}
% NAMESTRIP: no \namedateperiod here. The student writes their name once, on the
% cover this component is stapled behind. See references/conventions.md.

\vspace{0.15in}

\begin{notesbox}{Solving by Factoring --- On Your Own}
\small
Notes closed. Three items. Standard form first, factor completely, and list \emph{every} solution.

\begin{enumerate}[leftmargin=1.6em, itemsep=8pt, topsep=6pt]

\item Solve: $x^2-3x-10=0$
\begin{work}
  (x-5)(x+2) &= 0 \\
  x &= 5 \quad \text{or} \quad x = -2
\end{work}

\item Solve: $x^3-16x=0$. Then say how many solutions the equation has and what kind they are.
\begin{work}
  x\left(x^2-16\right) &= 0 \\
  x(x+4)(x-4) &= 0 \\
  x &= 0, \quad x = -4, \quad x = 4 \\
  \text{degree } 3 &\Rightarrow \text{ three solutions, all real}
\end{work}

\item A student solves $3x^2 = 12x$ by dividing both sides by $3x$, getting $x=4$. Is the answer
correct? Is it complete? Explain in one sentence, then give every solution.
\begin{work}
  3x^2-12x &= 0 \\
  3x(x-4) &= 0 \\
  x &= 0 \quad \text{or} \quad x = 4
\end{work}
\writeline

\end{enumerate}
\end{notesbox}

\end{document}
